\documentclass[10pt,a4paper]{amsart}
\usepackage[margin=19mm]{geometry}
\usepackage{amsmath,amssymb,mathtools}
\usepackage[scr=rsfs,cal=euler]{mathalfa}
\usepackage{xfrac}
\usepackage[unicode,hidelinks,bookmarks=false]{hyperref}
\newcommand{\ie}{i.e.\ }
\newcommand{\cf}{cf.\ }
\newcommand{\resp}{respectively\ }
\newcommand{\Subsection}{\subsection}
\providecommand{\nobreakdash}{\nobreak\hbox{-}}
\setlength{\emergencystretch}{2em}
\multlinegap0pt
\usepackage{amssymb}
\usepackage{enumerate}
\usepackage{tikz}
\usepackage{float}
\usepackage{mathtools}
\newtheorem{proposition}{Proposition}
\title{Span capacities of graphs}
\author{Mateja Gra\v{s}i\v{c}, Christopher Mouron, Aljo\v{s}a \v{S}uba\v{s}i\'c, Andrej Taranenko, Tanja Vojkovi\'c}
\date{Source: arXiv:2605.16852v1 (CC BY 4.0)}
\begin{document}
\maketitle

\noindent\emph{Source and licence.} Span capacities of graphs. Version arXiv:2605.16852v1 (CC BY 4.0), distributed by arXiv under the Creative Commons Attribution 4.0 International licence, http://creativecommons.org/licenses/by/4.0/, as recorded in the arXiv licence metadata for this exact version and shown on the abstract page https://arxiv.org/abs/2605.16852v1. Full licence notice: \texttt{licenses/arxiv-2605.16852-CC-BY-4.0.txt}.\par\smallskip

\noindent\textit{Editorial scope.} Counted unit: the article's proposition prop:strong-1 (source0.tex lines 469--471). The article's complete original proof of it is delivered with it (source0.tex lines 473--496, one \texttt{proof} environment of 24 lines). No mathematics is written, repaired or re-derived: the statement and the proof are the article's own lines, in the article's own notation, and no retained line is substituted or re-flowed.  No further source-local context is needed: every cross-reference of every retained line resolves inside the two delivered spans.  Nothing was omitted from any retained span, so the package carries no omission record.  No retained line contains a citation, so the wrapper carries no bibliography and no bibliography entry.  No retained line contains a figure environment, an image-inclusion command or drawing code, so no image is delivered and the package declares no asset dependency.  Editorial formatting only, in addition: an AMS \texttt{amsart} preamble that restates the article's own declarations and macros used by the retained lines, namely the environments \texttt{proposition} on separate counters, together with the standard packages those lines need.  The retained environments print in this document as Proposition 1 in order of appearance, whereas the article prints the same items with the numbering of its own full document.  The complete native source of the article as distributed in the arXiv e-print for this exact version is bundled unchanged as \texttt{sources/200750/source0.tex}.
\begin{proposition}\label{prop:strong-1}
    Any connected graph $G$ is strong-topfull.
\end{proposition}

\begin{proof}
   Suppose $G$ is a connected graph on $n$ vertices, say $v_1,\ldots,v_n$. We need to show that we can construct a collision-free weak $\ell$-tour $F=\{f_1,\ldots,f_n\}$ on $G$. 

   Let $f^1_i(0)=v_i$, for each $i\in\{1, 2, \ldots, n\}$, and let $f^1_1:\{0, 1, \ldots, \ell_1\}\rightarrow V(G)$ be an arbitrary $\ell_1$-track on $G$ such that $f^1_1(0)=v_1$. Since $G$ is connected, such a track (i.e. a walk through all vertices) clearly exists. For each $t\in\{1,\ldots,\ell_1\}$ and each $j\in\{2,\ldots,n\}$ define
   \[
    f^1_j(t) = \begin{cases}
                f^1_1(t-1), & \text{if } f^1_1(t) = f^1_j(t-1) \\
                f^1_j(t-1), & \text{otherwise}
             \end{cases}\text{.}
   \]
    It follows that $F^1=\{f^1_1, \ldots, f^1_n\}$ is a collision-free weak $\ell_1$-march on $G$ such that $f^1_1$ is an $\ell_1$-track.    

    We define the rest inductively. Assume $1<k\leq n$ and that $F^{k-1}$ is a collision-free weak $\ell_{k-1}$-march on $G$ such that $f^{k-1}_{k-1}$ is an $\ell_{k-1}$-track on $G$.
    Let $f^k_k:\{0, 1, \ldots, \ell_k\}\rightarrow V(G)$ be an arbitrary $\ell_k$-track on $G$ such that $f^k_k(0)=f^{k-1}_k(\ell_{k-1})$, and let $f^k_i(0)=f^{k-1}_i(\ell_{k-1})$, for all $i\in\{1,\ldots,n\}$. For each $t\in\{1,\ldots,\ell_k\}$ and each $j\in\{1,\ldots,k-1,k+1,\ldots,n\}$ define
   \[
    f^k_j(t) = \begin{cases}
                f^k_k(t-1), & \text{if } f^k_k(t) = f^k_j(t-1) \\
                f^k_j(t-1), & \text{otherwise}
             \end{cases}\text{.}
   \]
   Again, $F^k=\{f^k_1, \ldots, f^k_n\}$ is a collision-free weak $\ell_k$-march on $G$ such that $f^k_k$ is an $\ell_k$-track.

   It follows that $F:=F^n F^{n-1} \ldots F^1 = \{f_1, f_2, \ldots, f_n\}$ is a collision-free weak $\ell$-march on $G$, where $\ell=\ell_1+\ell_2+\ldots +\ell_n$. Moreover, $f_i$ is a track on $G$, for each $i\in \{1,\ldots,n\}$. Hence, $F$ is a collision-free weak $\ell$-tour on $G$.
\end{proof}

\end{document}
